\documentclass[11pt,a4paper]{article}
\usepackage[margin=27mm]{geometry}
\usepackage[T1]{fontenc}
\usepackage{lmodern,amsmath,amssymb,amsthm,mathtools,microtype}
\usepackage[dvipsnames]{xcolor}
\usepackage[colorlinks=true,linkcolor=MidnightBlue,urlcolor=MidnightBlue,citecolor=MidnightBlue]{hyperref}
\hypersetup{
  pdftitle={A limit-periodic Schrodinger operator with mixed absolutely continuous and singular continuous spectrum},
  pdfauthor={Matthew J. Colbrook},
  pdfsubject={Counterexample to AIM 1021: limit-periodic spectral purity}
}
\numberwithin{equation}{section}
\newtheorem{theorem}{Theorem}
\newtheorem{lemma}[theorem]{Lemma}
\theoremstyle{remark}
\newtheorem{remark}[theorem]{Remark}
\newcommand{\R}{\mathbb R}
\newcommand{\Z}{\mathbb Z}
\newcommand{\Id}{\mathrm I}
\newcommand{\tr}{\operatorname{tr}}
\newcommand{\dist}{\operatorname{dist}}
\newcommand{\spec}{\sigma}
\newcommand{\Leb}{\operatorname{Leb}}
\newcommand{\SL}{\mathrm{SL}}
\newcommand{\Ad}{\operatorname{Ad}}
\newcommand{\dd}{\,d}
\newcommand{\norm}[1]{\left\lVert #1\right\rVert}
\title{\textbf{A limit-periodic Schr\"odinger operator with\\
mixed absolutely continuous and singular continuous spectrum}}
\author{Matthew J. Colbrook\\[3pt]
{\small Department of Applied Mathematics and Theoretical Physics}\\
{\small University of Cambridge, United Kingdom}\\
{\small\href{mailto:mjc249@maths.cam.ac.uk}{\texttt{mjc249@maths.cam.ac.uk}}}}
\date{3 November 2026}
\begin{document}
\maketitle

\begin{abstract}
We describe a recursive construction of a bounded real potential on
$\Z$, uniformly approximated by periodic potentials, for which the
associated scalar discrete Schr\"odinger operator has nonzero absolutely
continuous and singular continuous spectral subspaces and no eigenvalues.
The mechanism combines adiabatic preservation of a local spectral density,
exponential narrowing of a different spectral cluster, and exact local
periodic repetitions. The quantitative point is that the protected
Floquet energies have a margin of order the inverse new period, whereas
the selected spectral bands have exponentially small widths.
\end{abstract}

\noindent\textbf{Problem and attribution.}
This construction gives a negative answer to
\href{https://github.com/MColbrook/AIM/blob/b014a51815ffd025d2b1abc4a3954205dfbf61e0/problems/1021-must-limit-periodic-discrete-schrodinger-operators-have-pure-spectral-type.md}{AIM problem 1021}
at repository revision \texttt{b014a51815ff}, namely Question~10.1 in
Damanik--Fillman~\cite{DF}. The construction is submitted here by
Matthew J. Colbrook. Periodic spectral formulas and the Gordon
argument are standard and are identified below.

\section{The construction theorem}
For a bounded real sequence $V$, let
\[
(H_Vu)(n)=u(n-1)+u(n+1)+V(n)u(n)
\]
on $\ell^2(\Z)$. Write $\mu_V$ for the spectral measure of $\delta_0$.
A sequence is limit-periodic if it is a uniform limit of periodic
sequences.

\begin{theorem}\label{thm:main}
There is a bounded real limit-periodic potential $V$ such that
\[
\mathcal H_{\mathrm{pp}}(H_V)=\{0\},\qquad
\mathcal H_{\mathrm{ac}}(H_V)\ne\{0\},\qquad
\mathcal H_{\mathrm{sc}}(H_V)\ne\{0\}.
\]
More specifically, there are a compact set $K$ of positive Lebesgue
measure and a constant $c>0$ such that
$\mu_V\ge c\mathbf1_K(E)\dd E$. There is also a bounded interval $I$
with endpoints in the resolvent set for which
$\mathbf1_I(H_V)\ne0$ and $\Leb(\spec(H_V)\cap I)=0$.
\end{theorem}

We construct periodic approximants $W_j$ of periods $p_j$, with
$p_j\mid p_{j+1}$. At each step a slow scalar offset is added to
successive $p_j$-site cells. One set of energies remains elliptic
throughout the slow cycle; its local spectral density is preserved
on a slightly smaller compact set. A second, isolated spectral
cluster encounters a long forbidden plateau at every energy; its
new bands are exponentially narrow. A zero plateau at the center
of the cycle preserves exact three-block repetitions.

\section{Periodic formulas}
For a real $p$-periodic $W$, put
\[
T_v(E)=\begin{pmatrix}E-v&-1\\1&0\end{pmatrix},\quad
M_W(E)=T_{W(p-1)}(E)\cdots T_{W(0)}(E),\quad
D_W(E)=\tr M_W(E).
\]
The monodromy transfers $(u(0),u(-1))^{\mathsf T}$ to
$(u(p),u(p-1))^{\mathsf T}$. Define the strictly elliptic set
\[
\mathcal E_W=\{E:|D_W(E)|<2\}.
\]
The spectrum is $\{|D_W|\le2\}$, a union of at most $p$ bands.
A closed Floquet gap may satisfy $|D_W|=2$ inside the interior
of the spectral set; thus strict ellipticity is the relevant
notion for the construction.

On $\mathcal E_W$, the density of $\mu_W$ is
\begin{equation}\label{eq:density}
\rho_W(E)=\frac{|(M_W(E))_{12}|}{\pi\sqrt{4-D_W(E)^2}}.
\end{equation}
To check the index and normalization, write
$M=\begin{psmallmatrix}a&b\\c&d\end{psmallmatrix}$. Bloch solutions
normalized by $u_\pm(0)=1$ satisfy
$u_\pm(-1)=(\lambda_\pm-a)/b$, where $\lambda_\pm$ are the
Floquet multipliers. Their Wronskian gives the diagonal Green
function, and its boundary imaginary part gives \eqref{eq:density}.
For $W=0$, this is $1/(\pi\sqrt{4-E^2})$; there is no extra
period factor. See also \cite[Chapter~7]{Teschl}.

Let
\[
J=\begin{pmatrix}0&-1\\1&0\end{pmatrix},\qquad R_\theta=e^{\theta J}.
\]
On each open component of $\mathcal E_W$ choose smooth real
$B(E)\in\SL(2,\R)$ and a real phase lift $\theta(E)$ with
\begin{equation}\label{eq:elliptic}
M_W(E)=B(E)R_{\theta(E)}B(E)^{-1}.
\end{equation}
There is no topological issue since the component is an interval.
For an explicit choice, set
\[
\tau=-\operatorname{sgn}(b)\sqrt{1-D_W^2/4},\quad
\beta=-b/\tau>0,\quad t=(a-d)/(2\tau),\quad
B=\begin{pmatrix}\sqrt\beta&0\\t/\sqrt\beta&1/\sqrt\beta\end{pmatrix},
\]
and take $\cos\theta=D_W/2$, $\sin\theta=\tau$.
The entry $b$ cannot vanish at a real elliptic matrix. In particular,
\begin{equation}\label{eq:beta}
(BJB^{-1})_{12}=-\beta,\qquad \rho_W=\beta/(2\pi).
\end{equation}

The signed phase has a strictly negative energy derivative.
Put $P_0=\Id$ and $P_r=T_{W(r-1)}\cdots T_{W(0)}$.
Differentiating the product gives
\[
M_W^{-1}(M_W)_E=-JC,\qquad
C=\sum_{r=0}^{p-1}P_r^{\mathsf T}e_1e_1^{\mathsf T}P_r.
\]
In the coordinates of \eqref{eq:elliptic}, the contribution from
differentiating $B$ is $R_\theta^{-1}LR_\theta-L$, where
$L=B^{-1}B_E$. Multiplication by $J$ and taking traces kills
that term. Consequently
\begin{equation}\label{eq:phase}
\theta_E=-\tfrac12\tr(B^{\mathsf T}CB)<0.
\end{equation}
The matrix $C$ is nonnegative and nonzero. Thus on compact
subsets of an elliptic component, $-\theta_E$ is bounded
above and below by positive constants.

\section{The slow profile}
Fix a smooth real $1$-periodic function $f$ with $|f|\le1$,
identically zero near $\Z$, and equal to $+1$ and $-1$
on two other intervals of positive length. Two disjoint
smooth bumps of opposite signs provide such a function.
All three plateau lengths are fixed.

For a $p$-periodic $W$, $\delta>0$, and an integer $N$, define
a $q=pN$-periodic potential by
\begin{equation}\label{eq:slow}
\mathcal S_{N,\delta}W(mp+r)=W(r)+\delta f(m/N),
\qquad 0\le r<p,\quad m\in\Z.
\end{equation}
Its distance from $W$ is at most $\delta$. The transfer matrix
across its $m$th old cell is exactly $M_W(E-\delta f(m/N))$.

\section{Adiabatic preservation of a bulk density}
Slow-deformation estimates of this kind appear in
Avila~\cite[Lemmas~2.11--2.12]{Avila}. We give the version
needed here, including a quantitative margin from every
new discriminant edge.

\begin{lemma}\label{lem:bulk}
Let $K$ be a finite union of compact intervals with
$K+[-\delta,\delta]\Subset\mathcal E_W$.
For every sufficiently small fixed $\eta>0$ and all sufficiently
large $N$, there is a finite union of compact intervals
$G_N\subset K$ such that, for $W_N=\mathcal S_{N,\delta}W$ and
$q=pN$,
\begin{align}
\Leb(K\setminus G_N)&\le C\eta,\label{eq:loss}\\
\sup_{G_N}|\rho_{W_N}-\rho_W|&\le C/(N\eta^2),\label{eq:rho}\\
\dist(G_N,\R\setminus\mathcal E_{W_N})&\ge c\eta^2/q.
\label{eq:margin}
\end{align}
The positive constants may depend on $W,p,\delta,K,f$, but not
on $N$ or $\eta$. The threshold for $N$ may depend on $\eta$.
\end{lemma}

\begin{proof}
Choose a slightly enlarged neighborhood $U$ of $K$ with
$\overline U+[-\delta,\delta]\Subset\mathcal E_W$.
Every shifted component stays in one old elliptic component.
Write
\[
A(s,E)=M_W(E-\delta f(s))
=B(s,E)R_{\theta(s,E)}B(s,E)^{-1},
\]
where $B(s,E)=B_W(E-\delta f(s))$ is chosen as in
\eqref{eq:elliptic}. Then $B(1,E)=B(0,E)=B_W(E)$.
All required derivatives are bounded on the compact parameter set.

Set $h=N^{-1}$ and decompose $B^{-1}B_s=gJ+S$, where $S$
is symmetric and traceless. On that two-dimensional real
matrix space, $\Ad_{R_\theta}$ is a rotation through $2\theta$.
Since $|\sin\theta|$ is bounded away from zero, the equation
\[
(\Ad_{R_\theta}-\Id)X=S
\]
has a unique smooth bounded solution in that space.
With $P(s,E)=e^{hX(s,E)}$ and coordinates
$U_m=B(mh,E)P(mh,E)Z_m$, one old-cell step becomes
\begin{align}
Z_{m+1}
&=P(s+h)^{-1}B(s+h)^{-1}B(s)R_\theta P(s)Z_m\notag\\
&=\bigl(R_{\theta-hg}+O(h^2)\bigr)Z_m.
\label{eq:normal}
\end{align}
Indeed $B(s+h)^{-1}B(s)=\Id-h(gJ+S)+O(h^2)$,
and the first-order nonrotational term is
$\Ad_{R_\theta}X-X-S=0$. The remainder is uniform.

Products of at most $N$ factors in \eqref{eq:normal}
are uniformly bounded by $(1+Ch^2)^N$. Telescoping against
their rotation parts gives error $O(Nh^2)=O(N^{-1})$.
Define
\[
\Phi_N(E)=\sum_{m=0}^{N-1}
\bigl(\theta(m/N,E)-N^{-1}g(m/N,E)\bigr).
\]
The full monodromy is therefore
\begin{equation}\label{eq:full}
M_{W_N}(E)
=B_W(E)\bigl(R_{\Phi_N(E)}+O(N^{-1})\bigr)B_W(E)^{-1}.
\end{equation}
In fact $X(0,E)=X(1,E)=0$ because the profile is constant
near those endpoints, so the endpoint correction is trivial.
All partial transfers inside the $q$-site cell and their inverses
are bounded independently of $N$. The at most $p$ extra
one-site factors within an old cell do not change this.

By \eqref{eq:phase}, for all sufficiently large $N$,
\[
c_1N\le-\Phi_N'(E)\le C_1N
\]
on each component of $U$. Set
$G_N=\{E\in K:|\sin\Phi_N(E)|\ge\eta\}$.
The monotone phase change of variables proves \eqref{eq:loss}:
in a phase interval of length $L$, the set $|\sin\phi|<\eta$
has length at most $C\eta(L+1)$. There are only finitely many
crossings, so $G_N$ is a finite union of compact intervals.

With $\beta=-(B_WJB_W^{-1})_{12}>0$, \eqref{eq:full} gives
\[
D_{W_N}=2\cos\Phi_N+O(N^{-1}),\qquad
(M_{W_N})_{12}=-\beta\sin\Phi_N+O(N^{-1}).
\]
On $G_N$, when $N^{-1}\ll\eta^2$, the trace lies strictly
between $-2$ and $2$. With $s=|\sin\Phi_N|\ge\eta$,
\[
\sqrt{4-D_{W_N}^2}=2s+O((Ns)^{-1}).
\]
Substitution in \eqref{eq:density}
gives $\rho_{W_N}=\beta/(2\pi)+O((N\eta^2)^{-1})$,
proving \eqref{eq:rho}.

Differentiate the $q$-site transfer product. It is a sum of
$q$ terms, each bounded uniformly by the bounds on left and
right partial transfers. Thus $|D_{W_N}'|\le Cq$ on $U$.
On $G_N$, the trace has margin $2-|D_{W_N}|\ge c\eta^2$.
For each $E\in G_N$, the interval of radius $c'\eta^2/q$
about $E$ lies in its component of $U$ for all large $N$.
Choosing $c'>0$ small enough, the derivative bound shows
that this whole interval remains strictly elliptic.
This proves \eqref{eq:margin}, including when $U$ is disconnected.
\end{proof}

\begin{remark}
The margin is measured from every point with $|D_{W_N}|\ge2$,
including closed gaps. Only the retained cores are used.
No convergence assertion across the moving band edges,
and no bound on spectral mass in moving exceptional sets,
is required.
\end{remark}

\section{Exponentially narrow target bands}
\begin{lemma}\label{lem:barrier}
Let $I$ be a compact interval and suppose for some $d>0$ that
\begin{equation}\label{eq:forbidden}
\max\{\dist(E,\spec(H_{W+\delta})),
\dist(E,\spec(H_{W-\delta}))\}\ge d
\quad (E\in I).
\end{equation}
Every portion in $I$ of a Floquet band of
$W_N=\mathcal S_{N,\delta}W$ has length at most $Ce^{-cN}$.
Consequently
\[
\Leb(\spec(H_{W_N})\cap I)\le C(pN)e^{-cN}.
\]
Constants are independent of $N$. If the endpoints of $I$
are in spectral gaps, the portions are complete bands.
\end{lemma}

\begin{proof}
Each $\pm1$ plateau contains a consecutive run of old cells
of total length $L\ge a q$, where $q=pN$ and $a>0$ is fixed,
for large $N$. On it the potential is exactly $W+\delta$
or $W-\delta$.

For a Floquet energy $E=E_k(\vartheta)\in I$,
$0<\vartheta<\pi$, normalize the Bloch solution by
\[
u(n+q)=e^{i\vartheta}u(n),\qquad
\sum_{n=0}^{q-1}|u(n)|^2=1.
\]
Then $|u(n)|\le1$ for all $n$. Choose the plateau operator
$H_U$ whose spectrum is at least $d$ from $E$. Cutting $u$
off to that plateau creates at most four boundary sources
for $H_U-E$, each of modulus at most one.

The elementary Combes--Thomas bound reads
\[
|\langle\delta_m,(H_U-E)^{-1}\delta_n\rangle|
\le\frac2d e^{-\gamma|m-n|},
\qquad \gamma=\operatorname{arsinh}(d/4).
\]
To see it, conjugate by a bounded clipped exponential weight:
the two hopping terms change in norm by at most
$2\sinh\gamma=d/2$, and the conjugated resolvent has norm
at most $2/d$. Thus at two adjacent middle sites of the
plateau,
\[
|u(r)|,\ |u(r+1)|\le\frac8d e^{-\gamma(L/2-2)}.
\]

The current $2\operatorname{Im}(u(n+1)\overline{u(n)})$
is independent of $n$. Hellmann--Feynman differentiation
of the boundary twist identifies it, up to sign, with
$E_k'(\vartheta)$ under the stated normalization. Hence
\[
|E_k'(\vartheta)|
\le128d^{-2}e^{4\gamma}e^{-\gamma L}.
\]
Scalar Floquet branches are simple and monotone on $(0,\pi)$.
Simplicity follows from the distinct multipliers
$e^{\pm i\vartheta}$. The current cannot vanish there,
since that would make the Bloch solution and its conjugate
linearly dependent. Integrating over $[0,\pi]$ proves the
band-portion estimate. There are at most $q$ bands.
\end{proof}

\begin{lemma}\label{lem:shift}
If $S$ is a union of at most $p$ compact intervals and
$r>p\Leb(S)$, there is $\delta\in(r,2r)$ such that
$(S+\delta)\cap(S-\delta)=\varnothing$.
\end{lemma}
\begin{proof}
If the component lengths are $\ell_1,\ldots,\ell_m$, then
\[
\Leb(S-S)\le\sum_{i,k}(\ell_i+\ell_k)
=2m\Leb(S)\le2p\Leb(S)<2r.
\]
Therefore $(2r,4r)$ contains a point outside the compact set
$S-S$. Take that point to be $2\delta$.
\end{proof}

Suppose $I$ isolates $S=\spec(H_W)\cap I$ from the rest of
the spectrum and $\delta$ is smaller than the isolation
margin. Disjointness of $S+\delta$ and $S-\delta$ then
implies \eqref{eq:forbidden} for some $d>0$: the two full
shifted spectra have no common point in $I$, so the
continuous maximum of the two distances has a positive
minimum on compact $I$. The value of $d$ may depend on
the chosen shift.

\section{Concrete initial data}
Take the period-three background
\[
W_0=(100,0,0),\qquad p_0=3,\qquad
D_0(E)=E^3-100E^2-3E+100.
\]
Fix
\[
I=[99,101],\qquad K_0=[0.989,0.991],\qquad
\delta_1=10^{-3}.
\]
The factorizations
\[
D_0(E)+2=(E-1)(E^2-99E-102),\qquad
D_0(E)-2=(E+1)(E^2-101E+98)
\]
show that the upper band $S_0$ lies in $(100,100.1)$,
while the other two bands lie in $[-2,2]$.
Since $D_0'\ge9997$ on the upper interval,
\[
\Leb(S_0)\le4/9997<0.000401<2\delta_1.
\]
The two translates $S_0\pm\delta_1$ are disjoint, giving
the barrier hypothesis on $I$.

For $E\in K_0$ and $|t|\le\delta_1$, one has
$E-t\in[0.988,0.992]$. On this interval $D_0'<0$, and
\[
D_0(0.988)=0.386030272,\qquad
D_0(0.992)=-0.406208512.
\]
Hence $K_0+[-\delta_1,\delta_1]\Subset\mathcal E_{W_0}$.
Let $m_0=\min_{K_0}\rho_{W_0}>0$.

Every later amplitude will be at most one quarter of
the preceding one. Therefore the total norm change from
$W_0$ is at most $(4/3)10^{-3}<0.002$.
The endpoints of $I$ stay in fixed spectral gaps,
the upper cluster remains nonempty, and the other
spectrum stays far below $I$.

\section{Induction with explicit parameter order}
Construct $W_j,p_j,K_j,\delta_{j+1}$ and write
\[
S_j=\spec(H_{W_j})\cap I,\qquad
g_j=\dist(K_j,\R\setminus\mathcal E_{W_j}).
\]
Fix positive tolerances $\alpha_j,\epsilon_j$ with
\[
\sum_{j\ge1}\alpha_j<\tfrac12\Leb(K_0),\qquad
\sum_{j\ge1}\epsilon_j<\tfrac12m_0.
\]
At the start of stage $j$, the data
$W_{j-1},p_{j-1},K_{j-1},\delta_j$ are fixed.
The shifted bulk is strictly elliptic and the shifted
target clusters are disjoint. This holds at $j=1$ by
the explicit initial data.

For every candidate $W_N=\mathcal S_{N,\delta_j}W_{j-1}$,
the already chosen amplitudes give
\[
\norm{W_N-W_0}_\infty\le\sum_{k=1}^j\delta_k
\le\tfrac43\delta_1<0.002.
\]
The same bound holds for $W_{j-1}\pm\delta_j$.
Their spectra are contained in
$[-2.002,2.002]\cup[99.998,100.102]$.
This makes the fixed isolation of $I$ a finite-stage invariant.

Apply Lemma~\ref{lem:bulk}. Choose $\eta_j>0$ first so
that the removed measure is less than $\alpha_j$.
Let $N$ tend to infinity, keeping those data and
$\eta_j$ fixed. The candidate outputs
\[
W_N=\mathcal S_{N,\delta_j}W_{j-1},\quad
p_N=p_{j-1}N,\quad K_N=G_N
\]
satisfy, for all large $N$,
\begin{align}
\Leb(K_{j-1}\setminus K_N)&<\alpha_j,\notag\\
\sup_{K_N}|\rho_{W_N}-\rho_{W_{j-1}}|&<\epsilon_j,\notag\\
g_N:=\dist(K_N,\R\setminus\mathcal E_{W_N})
&\ge\kappa_j/p_N,\label{eq:gN}\\
\Leb(S_N)&\le C_jp_Ne^{-c_jN}.\label{eq:sN}
\end{align}
The constants $\kappa_j,C_j,c_j>0$ are independent of
the new $N$.

For each candidate define
\begin{equation}\label{eq:radius}
r_N=\min\{g_N/8,\ 2^{-j-1}/p_N,\ \delta_j/8\}.
\end{equation}
Since $p_N\ge1$, \eqref{eq:gN} gives
\[
r_N\ge b_j/p_N,\qquad
b_j=\min\{\kappa_j/8,2^{-j-1},\delta_j/8\}>0.
\]
Meanwhile $p_N\Leb(S_N)\le C_jp_N^2e^{-c_jN}$.
Because $p_N=p_{j-1}N$, choose $N=N_j$ large enough
that
\begin{equation}\label{eq:compatibility}
p_N\Leb(S_N)<r_N.
\end{equation}
This is exactly the comparison $N^3e^{-c_jN}\to0$.
Also require each $N_j$ to exceed a fixed $N_*$, chosen
below for the repetitions.

Set $W_j=W_{N_j}$, $p_j=p_{N_j}$, $K_j=K_{N_j}$,
and $r_j=r_{N_j}$. Lemma~\ref{lem:shift} gives
$\delta_{j+1}\in(r_j,2r_j)$ separating the target
translates. Its upper bounds are
\begin{equation}\label{eq:deltabounds}
\delta_{j+1}<g_j/4,\qquad
\delta_{j+1}<2^{-j}/p_j,\qquad
\delta_{j+1}<\delta_j/4.
\end{equation}
The first ensures
$K_j+[-\delta_{j+1},\delta_{j+1}]\Subset\mathcal E_{W_j}$.
The fixed target isolation ensures the barrier hypothesis.
All constants for the following stage are now fixed
before $N_{j+1}$ is chosen. There is no circular
choice of parameters.

Since $\norm{W_j-W_{j-1}}_\infty\le\delta_j$, these
approximants converge uniformly to a real bounded $V$.
Their periods are nested and tend to infinity.
Thus $V$ is limit-periodic and $\norm{V}_\infty<100.002$.
Equivalently, the resulting potential is
\[
V(n)=W_0(n)+\sum_{j\ge1}\delta_j
f\left(\frac{\lfloor n/p_{j-1}\rfloor}{N_j}\right).
\]

\section{Surviving absolutely continuous spectrum}
The compact set $K_\infty=\bigcap_{j\ge0}K_j$ has
positive Lebesgue measure, by the summable losses.
For $E\in K_\infty$,
\[
\rho_{W_j}(E)\ge m_0-\sum_{k=1}^j\epsilon_k>m_0/2.
\]
Thus every approximating spectral measure dominates
the same nonzero absolutely continuous measure
\[
\nu=(m_0/2)\mathbf1_{K_\infty}(E)\dd E.
\]
Uniform convergence of potentials gives norm convergence
of their operators and weak convergence
$\mu_{W_j}\Rightarrow\mu_V$. Pass to the limit against
each nonnegative continuous test function in
$\mu_{W_j}\ge\nu$ to obtain $\mu_V\ge\nu$.
Consequently $\mathcal H_{\mathrm{ac}}(H_V)\ne\{0\}$.

\begin{remark}
Nonzero absolutely continuous spectrum has an independent
verification through reflectionless compactness. Each periodic
$W_j$ is reflectionless on the same set $K_\infty$.
The bounded reflectionless class on this fixed set is closed
under coefficientwise convergence; hence $V$ is reflectionless
there. Remling~\cite[Proposition~4.1(d)]{Remling} gives this compactness
on the fixed set $K_\infty$, with the common coefficient bound.
Part~(b) gives absolutely continuous spectrum for the half-line
restrictions: their Weyl functions have positive boundary imaginary
parts almost everywhere on $K_\infty$. The formula for the full-line
diagonal Green function in terms of these Weyl functions then gives a
nonzero absolutely continuous component for $H_V$; see also
\cite[Chapter~2]{Teschl}. This optional alternative is not used in the
direct measure proof.
The direct measure argument above additionally supplies the
stated lower bound on the chosen local spectral measure.
\end{remark}

\section{Surviving singular spectrum}
Put $t_j=\norm{V-W_j}_\infty$. Geometric decay and
\eqref{eq:deltabounds} give
\[
t_j\le\tfrac43\delta_{j+1},\qquad
p_jt_j\le\tfrac43\,2^{-j}\longrightarrow0.
\]
Hausdorff spectral continuity, together with the fixed
gaps at the endpoints of $I$, gives for large $j$
\[
\spec(H_V)\cap I\subset(S_j)_{t_j},
\]
where $(S_j)_{t_j}$ is the closed $t_j$-neighborhood.
There are at most $p_j$ intervals in $S_j$, so
\[
\Leb(\spec(H_V)\cap I)
\le\Leb(S_j)+2p_jt_j\longrightarrow0.
\]
Here $\Leb(S_j)\to0$ also follows directly from
\eqref{eq:compatibility}--\eqref{eq:radius}, which give
$\Leb(S_j)<2^{-j-1}/p_j^2$.

The cluster remains nonempty: the initial spectrum has
a point in $(100,100.1)$ and the total norm change is
less than $0.002$. Equivalently its isolated nonzero
spectral projection persists under this perturbation.
Therefore $\mathbf1_I(H_V)\ne0$, and its range is
entirely singular because it is supported on a null set.

In fact the same vector $\delta_0$ sees both spectral types.
Every modulation vanishes at the origin, so $V(0)=100$ and
\[
\int (E-100)^2\dd\mu_V(E)
=\norm{(H_V-100)\delta_0}^2=2.
\]
The spectrum outside $I$ lies in $[-2.002,2.002]$.
Consequently
\[
\mu_V(I)\ge1-\frac{2}{97.998^2}>0.99979.
\]
Thus $\mu_V$ has positive mass on the null spectral set
inside $I$, in addition to its absolutely continuous
component on $K_\infty$.

\section{Exact repetitions exclude eigenvalues}
Choose $a_0>0$ so that $f(s)=0$ whenever
$\dist(s,\Z)\le a_0$. Take $N_*\ge3$ with
$1/N_*<a_0$.

Fix $j$. On $[-p_j,2p_j-1]$ the old-cell indices
in the next update are $-1,0,1$, so that update
vanishes there. Every later update also vanishes
there: its cell length is at least $3p_j$, and
the indices are only $-1,0$. Hence
\begin{equation}\label{eq:exact}
V(n)=W_j(n)\quad(-p_j\le n\le2p_j-1).
\end{equation}

For a solution $H_Vu=Eu$, write
$U(n)=(u(n),u(n-1))^{\mathsf T}$ and let $A_j$
transfer across the middle block. Then
\[
U(p_j)=A_jU(0),\quad U(2p_j)=A_j^2U(0),\quad
U(-p_j)=A_j^{-1}U(0).
\]
Cayley--Hamilton gives
\[
A_j^2-(\tr A_j)A_j+\Id=0,\qquad
A_j+A_j^{-1}=(\tr A_j)\Id.
\]
If $|\tr A_j|\le1$, use the first identity; otherwise
use the second. In either case,
\[
\max\{\norm{U(-p_j)},\norm{U(p_j)},\norm{U(2p_j)}\}
\ge\tfrac12\norm{U(0)}.
\]
Since $p_j\to\infty$, no nonzero solution is square
summable. This is the exact three-block Gordon
argument; no rapid \emph{uniform} approximation rate
is assumed. See also \cite[Section~6.2]{DF}.

Thus $\mathcal H_{\mathrm{pp}}(H_V)=\{0\}$.
The nonzero singular projection found above is
therefore singular continuous. Together with the
absolutely continuous component, this proves
Theorem~\ref{thm:main}.

\section*{Scope}
The construction supplies a nonzero ac submeasure on
a compact set of positive measure, and a nonzero
purely sc restriction to the isolated interval $I$.
It does not classify every energy outside $I$ or
assert that the protected set contains an interval.
Neither assertion is needed for failure of purity.

\begin{thebibliography}{4}
\raggedright
\bibitem{DF}
D. Damanik and J. Fillman,
\emph{Spectral Properties of Limit-Periodic Operators},
arXiv:1802.05794 (2018); Chapter~15 in
\emph{Analysis and Geometry on Graphs and Manifolds},
London Mathematical Society Lecture Note Series~461,
Cambridge University Press (2020), pp.~382--444.
\href{https://doi.org/10.1017/9781108615259.016}{doi:10.1017/9781108615259.016}.
\url{https://arxiv.org/abs/1802.05794v1}.
\bibitem{Teschl}
G. Teschl,
\emph{Jacobi Operators and Completely Integrable Nonlinear Lattices},
Mathematical Surveys and Monographs 72,
American Mathematical Society (2000).
Author-hosted edition:
\url{https://www.mat.univie.ac.at/~gerald/ftp/book-jac/jacop.pdf}.
\bibitem{Avila}
A. Avila,
\emph{On the Kotani--Last and Schr\"odinger conjectures},
arXiv:1210.6325v1 (2012), Section~2.6, Lemmas~2.11--2.12.
\url{https://arxiv.org/abs/1210.6325v1}.
\bibitem{Remling}
C. Remling,
\emph{The absolutely continuous spectrum of Jacobi matrices},
arXiv:0706.1101, version~2 (2010), Proposition~4.1.
\url{https://arxiv.org/abs/0706.1101v2}.
\end{thebibliography}
\end{document}
